\section{Kết quả tối ưu mô hình}

\subsection{Thực nghiệm sơ bộ trên COCO128}
\label{sec:pilot-results}

Các thực nghiệm sơ bộ có mục đích kiểm chứng rằng từng giai đoạn của pipeline
hoạt động đúng và tìm cấu hình hợp lý trước khi chạy tốn kém trên tập PPE. Do
chạy trên CPU với tập rất nhỏ (COCO128 dùng chung cho train và validation), giá
trị tuyệt đối của $mAP$ không đại diện cho bài toán PPE; kết quả được dùng để so
sánh tương đối giữa các phương án. Tất cả số liệu dưới đây được đọc từ tệp
\texttt{results.csv} của các run, lấy epoch có $mAP@50{:}95$ cao nhất.

\textbf{Thưa hóa BatchNorm.} Bảng~\ref{tab:pilot-sparsity} cho thấy với lịch
$s$ giảm dần, mô hình vẫn duy trì độ chính xác cao trong khi $\gamma$ bị ép về 0;
đây là điều kiện để cắt tỉa không phá hỏng mô hình.

\begin{table}[H]
\centering
\caption{Thực nghiệm thưa hóa BatchNorm trên COCO128 (CPU)}
\label{tab:pilot-sparsity}
\resizebox{\textwidth}{!}{%
\begin{tabular}{lcccccc}
\toprule
\textbf{Cấu hình} & \textbf{Ảnh} & \textbf{Epoch} & \textbf{P} & \textbf{R} & \textbf{mAP@50} & \textbf{mAP@50:95}\\
\midrule
Sparsity $s=10^{-2}$, lịch giảm dần & 640 & 94/200 (tốt nhất 87) & 0,752 & 0,691 & 0,781 & 0,552\\
Sparsity $s=5\cdot10^{-2}$, SGD & 320 & 30 & 0,714 & 0,414 & 0,531 & 0,359\\
\bottomrule
\end{tabular}}
\end{table}

\textbf{Phục hồi sau cắt tỉa 30\%.} Bảng~\ref{tab:pilot-recovery} so sánh các
phương án phục hồi mô hình đã cắt tỉa (ảnh 320, AdamW, $lr_0=10^{-3}$). Hai điểm
đáng chú ý: (1) huấn luyện phục hồi ngắn (10 epoch) gần như không phục hồi được
($mAP \approx 0$), cho thấy mô hình sau cắt cần ngân sách huấn luyện đủ dài;
(2) với 100 epoch fine-tune, mô hình đã cắt đạt $mAP@50{:}95 = 0{,}383$. Các run
chưng cất dài bị dừng sớm do giới hạn thời gian CPU (mỗi epoch chưng cất tốn gấp
khoảng 2 lần fine-tune vì phải chạy thêm teacher), nên chưa thể kết luận chưng cất
kém hơn; ở cùng số epoch (khoảng 20), chưng cất với teacher đạt $0{,}185$.

\begin{table}[H]
\centering
\caption{Phục hồi độ chính xác sau cắt tỉa 30\% trên COCO128 (CPU, ảnh 320)}
\label{tab:pilot-recovery}
\resizebox{\textwidth}{!}{%
\begin{tabular}{lccccc}
\toprule
\textbf{Phương án} & \textbf{Epoch chạy} & \textbf{P} & \textbf{R} & \textbf{mAP@50} & \textbf{mAP@50:95}\\
\midrule
Fine-tune 10 epoch & 10 & 0,000 & 0,000 & 0,000 & 0,000\\
Chưng cất 10 epoch & 10 & 0,000 & 0,001 & 0,000 & 0,000\\
Fine-tune 100 epoch & 98 (tốt nhất 87) & 0,676 & 0,448 & 0,515 & 0,383\\
Chưng cất 100 epoch (dừng sớm) & 21 & 0,623 & 0,238 & 0,281 & 0,185\\
Chưng cất, teacher YOLO26x (dừng sớm) & 7 & 0,605 & 0,117 & 0,158 & 0,099\\
\bottomrule
\end{tabular}}
\end{table}

\textbf{Độ nhạy của trọng số chưng cất.} Để chọn $\lambda_{cwd}$, chưng cất được
chạy trên YOLO26n chưa cắt tỉa trong 5 epoch với nhiều giá trị $\lambda$
(Bảng~\ref{tab:pilot-lambda}). Độ chính xác giảm đơn điệu khi $\lambda$ tăng; với
$\lambda = 10$ $mAP@50{:}95$ chỉ còn 0,075. Trong ngân sách huấn luyện ngắn, số
hạng CWD lớn lấn át mất mát phát hiện, nên $\lambda$ nhỏ phù hợp hơn cho giai đoạn
đầu; giá trị $\lambda=50$ mặc định chỉ phù hợp khi huấn luyện đủ dài và nên được
dò lại trên tập PPE.

\begin{table}[H]
\centering
\caption{Ảnh hưởng của $\lambda_{cwd}$ (YOLO26n, 5 epoch, COCO128, ảnh 320)}
\label{tab:pilot-lambda}
\begin{tabular}{lcc}
\toprule
\textbf{$\lambda_{cwd}$} & \textbf{mAP@50} & \textbf{mAP@50:95}\\
\midrule
0,5 & 0,384 & 0,252\\
0,5 (teacher YOLO26x) & 0,386 & 0,255\\
1,0 & 0,366 & 0,227\\
2,0 & 0,362 & 0,209\\
5,0 & 0,307 & 0,188\\
10,0 & 0,141 & 0,075\\
\bottomrule
\end{tabular}
\end{table}

\textbf{Độ biến thiên theo seed.} Chưng cất ($\lambda$ nhỏ) và fine-tune thường
được lặp với 3 seed (0, 1, 2), 5 epoch (Bảng~\ref{tab:pilot-seed}). Chênh lệch
trung bình giữa hai phương án ($\approx 0{,}01$ $mAP@50{:}95$) nhỏ hơn độ lệch giữa
các seed của cùng một phương án (tới $0{,}04$). Vì vậy so sánh giữa các cấu hình
tối ưu phải dựa trên nhiều seed hoặc huấn luyện đủ dài; kết luận từ một run ngắn
duy nhất là không tin cậy.

\begin{table}[H]
\centering
\caption{Độ biến thiên theo seed (YOLO26n, 5 epoch, COCO128)}
\label{tab:pilot-seed}
\begin{tabular}{lccccc}
\toprule
\textbf{Phương án} & \textbf{Seed 0} & \textbf{Seed 1} & \textbf{Seed 2} & \textbf{Trung bình} & \textbf{Chỉ số}\\
\midrule
Chưng cất & 0,212 & 0,247 & 0,209 & 0,223 & mAP@50:95\\
Fine-tune & 0,248 & 0,240 & 0,211 & 0,233 & mAP@50:95\\
Chưng cất & 0,345 & 0,405 & 0,357 & 0,369 & mAP@50\\
Fine-tune & 0,370 & 0,367 & 0,355 & 0,364 & mAP@50\\
\bottomrule
\end{tabular}
\end{table}

\subsection{Kết quả trên tập Construction-PPE}

Bảng~\ref{tab:model-accuracy-benchmark} tổng hợp kết quả của các cấu hình ở
Bảng~\ref{tab:optimization-configs} trên tập test Construction-PPE, đo bằng mô
hình PyTorch (fake-quant đối với cấu hình QAT). Số tham số và FLOPs là giá trị đo
được từ checkpoint.

\begin{table}[H]
\centering
\caption{Kết quả các cấu hình tối ưu trên tập test Construction-PPE}
\label{tab:model-accuracy-benchmark}
\resizebox{\textwidth}{!}{%
\begin{tabular}{lccccc}
\toprule
\textbf{Cấu hình} & \textbf{Params} & \textbf{GFLOPs} & \textbf{Precision} &
\textbf{mAP@50} & \textbf{mAP@50:95}\\
\midrule
Baseline FP32 & 2,57 M & 6,12 & \cbs{} & \cbs{} & \cbs{}\\
Pruned-30 + fine-tune & 2,17 M & 4,82 & \cbs{} & \cbs{} & \cbs{}\\
Pruned-30 + distill & 2,17 M & 4,82 & \cbs{} & \cbs{} & \cbs{}\\
Pruned-30 + distill + QAT & 2,17 M & 4,82 & \cbs{} & \cbs{} & \cbs{}\\
\bottomrule
\end{tabular}}
\end{table}
